% GENERATED FILE. Edit canonical lessons, then run npm run generate:books.
% canonical-source-hash: fnv1a64:b7d6e86d01b603e6
% canonical-lessons: AR-C19-wahid-khamsa

\chapter{Numbers One to Five}
\label{ch:ar-numbers-one-to-five}
\hlchaptermodality{19}

\begin{chapteropening}
\textbf{By the end of this chapter:} \emph{I can count from one to five in Arabic.}

Count wāḥid to khamsa, then separate the Arabic number WORDS from the Indian-origin digit shapes.
\end{chapteropening}

\section[\ar{واحد} \ar{اثنان} \ar{ثلاثة} \ar{أربعة} \ar{خمسة}]{\ar{واحد،} \ar{اثنان،} \ar{ثلاثة،} \ar{أربعة،} \ar{خمسة} — one through five}

\label{lesson:AR-C19-wahid-khamsa}



\begin{quote}
Before the numbers themselves — an honest correction to a very common mix-up: \textquotedblleft{}Arabic numerals\textquotedblright{} (the digit shapes \textbf{0, 1, 2, 3}...) are \textbf{not} these words, and don't come from Arabic at all.
\end{quote}

\subsection*{You'll want to know — \ar{واحد}, \ar{اثنان}, \ar{ثلاثة}, \ar{أربعة}, \ar{خمسة} — the words}
\noindent
\begin{tabularx}{\linewidth}{@{}>{\raggedright\arraybackslash}X>{\raggedright\arraybackslash}X>{\raggedright\arraybackslash}X@{}}
\toprule
\textbf{Arabic} & \textbf{Transliteration} & \textbf{Meaning} \\
\midrule
\textbf{\ar{واحد}} & \emph{wāḥid} & one \\
\textbf{\ar{اثنان}} & \emph{ithnān} & two \\
\textbf{\ar{ثلاثة}} & \emph{thalātha} & three \\
\textbf{\ar{أربعة}} & \emph{arbaʿa} & four \\
\textbf{\ar{خمسة}} & \emph{khamsa} & five \\
\bottomrule
\end{tabularx}

These are ordinary Semitic-root Arabic words, built the same way Arabic builds any noun — nothing to do with the shapes of the digits you write.

\begin{culture}
Here's the correction: the \textbf{digit shapes} \textbf{0, 1, 2, 3, 4, 5, 6, 7, 8, 9} that English calls \textquotedblleft{}Arabic numerals\textquotedblright{} didn't originate in Arabic at all — they were invented in \textbf{India}, appearing in Indian mathematical texts and inscriptions centuries before reaching the Arab world. The 9th-century Persian mathematician \textbf{al-Khwārizmī} (whose name gives English \textbf{algorithm}) wrote an influential treatise on calculating with these \textbf{Hindu numerals}, which carried the system into the Islamic world and then, via Latin translations, into Europe. So \textquotedblleft{}Arabic numerals\textquotedblright{} is really a transmission credit, not an origin credit — the shapes are Indian; only their route to Europe ran through Arabic mathematics. The \textbf{words} you're learning here (\emph{wāḥid, ithnān, thalātha}...) are a completely separate thing: genuine, native Arabic vocabulary, unrelated to that digit-shape history.
\end{culture}

\subsection*{Your turn}
Say these aloud:

\begin{itemize}
\raggedright
  \item \textquotedblleft{}wāḥid, ithnān, thalātha, arbaʿa, khamsa\textquotedblright{} — one through five
  \item the honest correction — \textquotedblleft{}Arabic numerals\textquotedblright{} (the digit SHAPES) are actually Indian in origin, not these Arabic WORDS
\end{itemize}

\begin{tcolorbox}[breakable,colback=teal!4,colframe=teal!35!black,title={Before you move on}]
What are the Arabic words for one through five? (\textbf{wāḥid, ithnān, thalātha, arbaʿa, khamsa}.) Where did the digit shapes 0-9 actually originate, and who's credited with carrying them to Europe? (\textbf{India} — carried into the Islamic world and then Europe largely through \textbf{al-Khwārizmī}'s writings.) Are \textquotedblleft{}Arabic numerals\textquotedblright{} the same thing as these Arabic number words? (\textbf{No} — completely separate: one is a set of written digit shapes with Indian origins, the other is native Arabic vocabulary.)
\end{tcolorbox}
